\section{问题重述}

\subsection{问题背景}

微网是为解决光伏发电等分布式能源本地消纳、减少并网问题而快速发展起来的小型发用电系统。出于经济性考虑，非孤岛式微网可以在外网电价较低或分布式能源有剩余电量的时段为储能充电，并在外网电价较高或光伏不足的时段释放电能。由此，微网与外网的调控目标可以概括为：利用小区负载、光伏发电量与储能当前储电量等数据，尽可能精准地给出微网的购电量和储能充放电量，使供电既满足小区负载，又尽可能节省购电费用。

本问题的物理对象由三部分组成：光伏发电单元、储能设备与小区负载，并通过公共连接点与外网相连。附录 1 给定储能参数：最大容量 $12000$ kWh，最大充放电功率 $5000$ kW，2025 年 1 月 1 日 0:00 储电量为 $6000$ kWh，运行期储电量必须保持在 $1200\sim10800$ kWh 之间，充放电效率均为 $90\%$。附件数据的时间分辨率均为 $10$ min，即一天可分为 $144$ 个调度时段。

问题的核心困难在于\textbf{信息的可获得性}：微网必须在决策时刻仅依据当时已知的信息制定策略，而负载、光伏与实际电价在一天内是逐步揭晓的。因此四个问题的差异并不在物理约束，而在于\emph{决策时刻能看到什么信息}，以及\emph{计划与实际不一致时如何结算}。

\subsection{问题提出}

\noindent \textbf{问题 1}：每天的电价和小区负载相同，并给出光伏发电功率一天的预测数据。微网提供的电能不可低于小区负载，且储能设备在 $0{:}00$ 和 $24{:}00$ 的储电量相同。要求在每天 $0{:}00$ 制定微网当天的计划购电策略，给出指定时段的购电量、全天购电量与购电费，以及储能的充放电量与首末储电量。

\noindent \textbf{问题 2}：每天的电价相同，而小区负载和光伏发电功率随时间变化。微网提供的电能不可低于小区负载；若低于负载，需向外网紧急购电，其电价是交易时刻电价的 $5$ 倍。除紧急购电费用外，其他时间段的购电费用均按计划购电量计算。要求在每天 $0{:}00$ 制定当天的计划购电策略，给出指定日期（2025.3.20、2025.6.21、2025.9.23、2025.12.21）的购电量、储能充放电量，并给出紧急购电结果。

\noindent \textbf{问题 3}：在每天 $0{:}00$、$6{:}00$、$12{:}00$ 和 $18{:}00$ 可获得未来 $24$ 小时整点的光伏发电功率预报。可根据 $0{:}00$ 的预报制定当天计划购电策略，也可根据其他时刻的预报调整购电策略。计划购电量高于调整购电量的部分，违约电价是交易时刻电价的 $50\%$；调整购电量高于计划购电量的部分，超出部分的电价是交易时刻电价的 $1.5$ 倍。总购电费用包括计划购电费用、紧急购电费用和调整购电量的相关费用。要求建立微网当天购电策略的数学模型，并\textbf{分析是否需要引入其他时刻的预报制定调整购电策略}。

\noindent \textbf{问题 4}：外网电价也实时波动。要求针对波动电价建立微网当天购电策略的数学模型，根据附件 4 的电价数据重新计算问题 2 和问题 3，并给出相应的计算结果。

\subsection{研究现状综述}

微网经济调度问题在数学上通常被建模为以购电费或综合运行成本最小为目标的线性规划或混合整数规划，外网购电、分布式电源消纳与储能充放电被统一纳入同一组能量平衡约束中 \cite{wu2014dynamic,nemati2018optimization}。确定性模型适用于负载、电价与光伏均已知的日前场景，其优点是约束结构清晰、求解稳定可靠。

当光伏出力存在预测误差时，主流处理方式有三类：一是\textbf{情景法／两阶段随机规划}，用有限个历史或生成情景刻画不确定性，第一阶段确定与情景无关的日前计划，第二阶段按情景确定执行量 \cite{wu2016solution,sun2025distributionally}；二是\textbf{鲁棒优化}，在最坏情景下保证可行性，代价是结果偏保守 \cite{zhang2018robust}；三是\textbf{滚动时域优化（模型预测控制）}，在新信息到达后仅重优化剩余时域并只执行近期决策，天然适配“预报分批发布”的决策结构 \cite{elkazaz2020energy,silva2020optimal}。含光伏与储能的日前调度研究普遍指出，储能的价值高度依赖预测精度与调度时域长度 \cite{luo2020optimal}，而储能的运行成本又与其循环深度密切相关 \cite{maheshwari2020optimizing}。

针对“预测值应偏向哪一侧”的问题，报童模型给出了经典答案：当高估与低估的边际损失不对称时，最优的点预测不是条件均值，而是某个\textbf{分位数} \cite{qin2011newsvendor}。这一结论在电力系统中已被用于发电商的投标策略与偏差电量定价 \cite{polat2024renewable}。在预测方法层面，基于数值天气预报的机器学习模型是目前光伏功率预测的主流 \cite{markovics2022comparison}；在市场机制层面，实时电价与需求响应的耦合会显著改变用户的购电行为 \cite{wang2015review}。

现有研究多聚焦于算法本身，而对\textbf{信息集边界}与\textbf{结算口径}的刻画往往不够严格：不少工作在日前优化中直接使用当天的真实负载或真实电价，使结果成为不可实现的“完美信息下界”；也有工作在比较不同策略时同时改变了预测方法、约束条件与费用口径，导致差异无法归因。本文针对这两点，把信息集边界作为建模的第一原则，并通过单因素消融实验保证每条结论都可归因。此外，与已有工作多采用元启发式求解不同，本文利用问题的线性结构精确求解，避免了启发式算法结果不可复现的问题。
